% !TEX program = xelatex
% 공통수학2 · Ⅲ. 함수와 그래프 · 1. 함수 · 중단원 마무리 (38차시)
% 학생용 활동지 (답안 없음)
\documentclass[11pt,a4paper]{article}
\usepackage{kotex}
\usepackage{iftex}
\ifPDFTeX\else
  \setmainhangulfont{Noto Serif CJK KR}
  \setsanshangulfont{Noto Sans CJK KR}
\fi
\usepackage[margin=2.1cm]{geometry}
\usepackage{parskip}
\usepackage{enumitem}
\usepackage{array}
\usepackage{tabularx}
\usepackage{xcolor}
\usepackage{amssymb}
\usepackage{tikz}
\usepackage[most]{tcolorbox}
\usepackage{fancyhdr}
\usepackage{titlesec}

\definecolor{goldc}{HTML}{B07C22}
\definecolor{goldstrong}{HTML}{8A5F16}
\definecolor{indigoc}{HTML}{2B3B57}
\definecolor{indigosoft}{HTML}{6E82A6}
\definecolor{linec}{HTML}{D6DACB}
\definecolor{surface2}{HTML}{E4E8DC}
\definecolor{writeline}{HTML}{C7CCBA}
\definecolor{inksoft}{HTML}{52604F}

\titleformat{\section}{\normalfont\sffamily\bfseries\large\color{indigoc}}{\thesection}{0.6em}{}
\titlespacing{\section}{0pt}{14pt}{6pt}
\newtcolorbox{goalbox}{enhanced, breakable, colback=white, colframe=goldc, boxrule=0.5pt, leftrule=3pt, arc=2pt, left=10pt, right=10pt, top=8pt, bottom=8pt}
\newtcolorbox{sheetbox}{enhanced, breakable, colback=white, colframe=linec, boxrule=0.6pt, arc=3pt, left=11pt, right=11pt, top=9pt, bottom=9pt}
\newcommand{\writespace}[1]{%
  \par\nobreak\vspace{3pt}
  \foreach \n in {1,...,#1} {%
    \noindent\textcolor{writeline}{\rule{\linewidth}{0.4pt}}\par\vspace{0.62cm}%
  }
}
\newcommand{\fillblank}[1]{\makebox[#1]{\hrulefill}}
\pagestyle{fancy}\fancyhf{}\renewcommand{\headrulewidth}{0pt}
\fancyfoot[L]{\footnotesize\color{inksoft} 공통수학2 · 2026학년도 2학기 · 지평선고등학교}
\fancyfoot[R]{\footnotesize\color{inksoft} 다음 시간 --- 2. 유리함수와 무리함수}

\begin{document}

\noindent
\begin{minipage}[b]{0.62\linewidth}
  {\sffamily\footnotesize\color{indigosoft} 공통수학2 · Ⅲ. 함수와 그래프}\\[2pt]
  {\sffamily\bfseries\Large 1. 함수 \textperiodcentered\ 중단원 마무리}
\end{minipage}\hfill
\begin{minipage}[b]{0.36\linewidth}
  \raggedleft\small\color{inksoft}
  반\ \fillblank{1.6cm} \quad 번호\ \fillblank{1.6cm}\\[4pt]
  이름\ \fillblank{3.6cm}
\end{minipage}

\vspace{4pt}
\noindent{\color{goldc}\rule{\linewidth}{2.2pt}}
\vspace{10pt}

\begin{goalbox}
\textbf{오늘의 목표}\\
33$\sim$37차시에서 배운 개념(함수의 뜻, 여러 가지 함수, 합성함수, 역함수)을 스스로 정리하고 점검한다.
\end{goalbox}

\section{개념 지도 채우기}

\begin{sheetbox}
\begin{tabularx}{\linewidth}{@{}>{\bfseries\small}p{4.6cm}X@{}}
\toprule
\textbf{개념} & \textbf{핵심 내용} \\
\midrule
함수의 뜻 & \fillblank{6.4cm} \\[6pt]
정의역·공역·치역 & \fillblank{6.4cm} \\[6pt]
일대일함수 / 일대일대응 & \fillblank{6.4cm} \\[6pt]
합성함수 $(g\circ f)(x)$ & \fillblank{6.4cm} \\[6pt]
역함수의 존재 조건과 표기 & \fillblank{6.4cm} \\[6pt]
역함수의 그래프 & \fillblank{6.4cm} \\
\bottomrule
\end{tabularx}
\end{sheetbox}

\section{기본 문제}

\begin{sheetbox}
\textbf{1.} $X=\{1,2,3,4\}$, $Y=\{1,2,\ldots,10\}$, $f(x)=2x$일 때, 정의역, 공역, 치역을 각각 구하시오.
\writespace{2}

\textbf{2.} 다음 중 일대일함수인 것을 모두 고르시오.\\
\hspace*{1.6em}① $f(x)=5-x$ \quad ② $f(x)=x^2-1$ \quad ③ $f(x)=|x|$ \quad ④ $f(x)=x^3$
\writespace{2}

\textbf{3.} $f(x)=3x+2$, $g(x)=-x+1$일 때, $(g\circ f)(1)$과 $(f\circ g)(1)$을 각각 구하시오.
\writespace{2}

\textbf{4.} $f(x)=4x-3$의 역함수 $f^{-1}(x)$를 구하시오.
\writespace{2}
\end{sheetbox}

\section{표준 \& 도전 문제}

\begin{sheetbox}
\textbf{5.} 증가함수 $f(x)=2x-6$의 그래프와 그 역함수의 그래프가 만나는 점을 구하시오.
\writespace{2}

\textbf{6. (서술형)} $f(x)=x+3$, $g(x)=2x$일 때, $(g\circ f)^{-1}(x)$를 직접 구한 값과 $f^{-1}\circ g^{-1}$을 이용해 구한 값이 같음을 보이시오.
\writespace{4}

\textbf{7.} $X=\{1,2,3\}$에서 항등함수와 상수함수(값 3)를 각각 순서쌍의 집합으로 나타내시오.
\writespace{2}

\textbf{8. (도전)} $f(x)=ax+b$ $(a\neq0)$가 자기 자신의 역함수, 즉 $f=f^{-1}$가 되도록 하는 $a$의 값을 구하시오.
\writespace{3}
\end{sheetbox}

\section{마무리 성찰 --- 나의 성취 수준 점검}

\begin{sheetbox}
아래 항목 중 자신 있게 설명할 수 있는 것에 $\checkmark$ 표시해 보자.

\begin{itemize}[leftmargin=1.6em, itemsep=3pt]
  \item[$\square$] 함수의 뜻(존재성·유일성)을 알고 정의역·공역·치역을 구별할 수 있다.
  \item[$\square$] 일대일함수, 일대일대응, 항등함수, 상수함수를 구별할 수 있다.
  \item[$\square$] 두 함수의 합성함수를 구할 수 있다.
  \item[$\square$] 역함수의 존재 조건을 알고, 역함수를 구할 수 있다.
  \item[$\square$] 함수와 역함수의 그래프가 $y=x$ 대칭임을 알고, 교점을 구할 수 있다.
\end{itemize}

\vspace{6pt}
{\footnotesize\color{inksoft} 가장 보충이 필요한 부분과 그 이유를 적어 보자.}
\writespace{2}
\end{sheetbox}

\end{document}
